\documentclass{book}
\usepackage{amsmath, amssymb, amsthm}
\usepackage{geometry}
\usepackage[CJKspace]{xeCJK}  % 한글 지원
\setmainfont{Times New Roman}
\setCJKmainfont{Malgun Gothic}

\geometry{
 a4paper,
 left=10mm,
 right=10mm,
 top=15mm,
 bottom=15mm
}
\setlength{\parindent}{0pt}

\title{FNN v2 \& Hyper-FNN (Matrix Form)}
\author{사서림}
\date{2025.08.27.}

\begin{document}
\maketitle
\tableofcontents

\chapter{FNN v2 - Fourier Feature Neural Network (Matrix Form)}

\section{기본 정의}
입력 $V \in \mathbb{R}^n$, 출력 $W \in \mathbb{R}^m$.
푸리에 급수 항의 개수를 $S$라 하자.

\[
\omega := \frac{2\pi}{C \cdot m}, \qquad
X = \text{Pad}(V; S)
\]

여기서 $\text{Pad}$는 길이를 $S$로 맞추기 위한 패딩 연산이다.

\section{Fourier Feature 벡터}
입력 $X \in \mathbb{R}^S$에 대하여 다음의 푸리에 특징 벡터를 정의한다.
\[
\Phi(X) =
\begin{bmatrix}
\sin(\Omega X + H_s) \\
\cos(\Omega X + H_c)
\end{bmatrix}
\in \mathbb{R}^{2S}
\]
여기서 $\Omega, H_s, H_c \in \mathbb{R}^{S \times S}$는 학습 파라미터 행렬이다.

\section{출력 수식}
FNN v2의 출력은 다음과 같이 정의된다.
\[
W = A \, \Phi(X) + b + \lVert V \rVert
\]
여기서 $A \in \mathbb{R}^{m \times 2S}$, $b \in \mathbb{R}^m$은 학습 파라미터이다.

\section{헤더 MLP}
성능 향상을 위해, $W$에 대해 2단계 MLP를 적용한다.
\[
M(W) = L_2 \circ \sigma \circ L_1(W),
\]
\[
L_1:\mathbb{R}^m \to \mathbb{R}^{8m}, \quad
L_2:\mathbb{R}^{8m} \to \mathbb{R}^m, \quad
\sigma=\text{GELU}.
\]

---

\section{Hyper-FNN v2 - Hypernetwork-based FNN (Matrix Form)}

\subsection{하이퍼네트워크의 정의}
하이퍼네트워크 $H$는 입력 $V \in \mathbb{R}^n$을 받아
FNN v2의 모든 파라미터 행렬을 동적으로 생성한다.
\[
H: V \mapsto \{A(V), \Omega(V), H_s(V), H_c(V), b(V)\}.
\]

\subsection{순전파}
생성된 파라미터를 이용해 출력은 다음과 같이 정의된다.
\[
W(V) = A(V) \, \Phi\!\big(X; \Omega(V), H_s(V), H_c(V)\big) + b(V) + \lVert V \rVert.
\]

\subsection{의의}
\begin{itemize}
\item \textbf{입력 적응성}: 파라미터가 $V$에 의존하므로 데이터 분포 변화에 강함.
\item \textbf{파라미터 효율성}: 거대한 파라미터를 직접 저장하지 않고, 작은 하이퍼네트워크 $H$만 학습하면 됨.
\end{itemize}

---

\end{document}
